\appendix

%================================================
\chapter{8086 Instruction Summary}
%================================================

This appendix puts in one place all the instructions we studied in Chapter 3, so you can find them quickly while
you write your programs. Symbols: d = destination, s = source, r = register, m = memory box.

\section*{The most used DOS functions (\code{int 21h})}
\begin{center}
\small
\begin{tabular}{|c|l|l|}
    \hline
    \ras{AH} & \ras{Function} & \ras{Inputs / outputs} \\ \hline
    \hex{01} & read a key and show it & output in AL \\ \hline
    \hex{02} & print a character & the character in DL \\ \hline
    \hex{09} & print a string ending with \code{\$} & the address in DX \\ \hline
    \hex{0A} & read a line into a buffer & the address in DX \\ \hline
    \hex{4C} & end the program & the exit code in AL \\ \hline
\end{tabular}
\end{center}

\section*{The instructions}
\begin{center}
\small\renewcommand{\arraystretch}{1.12}
\begin{tabularx}{\linewidth}{|l|X|l|}
    \hline
    \ras{Instruction} & \ras{Meaning} & \ras{Flags changed} \\ \hline
    \multicolumn{3}{|c|}{\cellcolor{brandOrange!20}\textbf{Transfer}} \\ \hline
    \code{mov d, s} & \nql{d}{s} & none \\ \hline
    \code{xchg a, b} & swap a and b & none \\ \hline
    \code{lea r, m} & r gets the address of m & none \\ \hline
    \code{push s} & \nql{SP}{SP - 2}, then \nql{[SP]}{s} & none \\ \hline
    \code{pop d} & \nql{d}{[SP]}, then \nql{SP}{SP + 2} & none \\ \hline
    \code{in al, port}, \code{out port, al} & read from a port, write to a port & none \\ \hline
    \multicolumn{3}{|c|}{\cellcolor{brandOrange!20}\textbf{Arithmetic}} \\ \hline
    \code{add d, s} & \nql{d}{d + s} & all \\ \hline
    \code{adc d, s} & \nql{d}{d + s + CF} & all \\ \hline
    \code{sub d, s} & \nql{d}{d - s} & all \\ \hline
    \code{sbb d, s} & \nql{d}{d - s - CF} & all \\ \hline
    \code{inc d} & \nql{d}{d + 1} & all except CF \\ \hline
    \code{dec d} & \nql{d}{d - 1} & all except CF \\ \hline
    \code{neg d} & \nql{d}{-d} & all \\ \hline
    \code{cmp a, b} & computes a - b without keeping it & all \\ \hline
    \code{mul s} (8 bits) & \nql{AX}{AL * s} & CF, OF \\ \hline
    \code{mul s} (16 bits) & \nql{DX:AX}{AX * s} & CF, OF \\ \hline
    \code{div s} (8 bits) & \nql{AL}{AX / s}, \nql{AH}{AX mod s} & undefined \\ \hline
    \code{div s} (16 bits) & \nql{AX}{DX:AX / s}, \nql{DX}{DX:AX mod s} & undefined \\ \hline
    \multicolumn{3}{|c|}{\cellcolor{brandOrange!20}\textbf{Logic and shifts}} \\ \hline
    \code{and}, \code{or}, \code{xor} & bit by bit operations & CF = OF = 0, ZF, SF \\ \hline
    \code{not d} & flip all the bits & none \\ \hline
    \code{test a, b} & \code{and} without keeping the result & like \code{and} \\ \hline
    \code{shl d, n}, \code{shr d, n} & shift: multiply or divide by $2^n$ & CF = last bit out \\ \hline
    \code{rol d, n}, \code{ror d, n} & rotate & CF \\ \hline
    \multicolumn{3}{|c|}{\cellcolor{brandOrange!20}\textbf{Flow control}} \\ \hline
    \code{jmp L} & always jump & none \\ \hline
    \code{je}, \code{jne}, \code{jz}, \code{jnz} & depends on ZF & -- \\ \hline
    \code{jg}, \code{jge}, \code{jl}, \code{jle} & \textbf{signed} comparison & -- \\ \hline
    \code{ja}, \code{jae}, \code{jb}, \code{jbe} & \textbf{unsigned} comparison & -- \\ \hline
    \code{loop L} & \nql{CX}{CX - 1}, then jump if CX $\neq 0$ & none \\ \hline
    \code{jcxz L} & jump if CX = 0 & -- \\ \hline
    \code{call P}, \code{ret} & call a procedure, come back from it & none \\ \hline
    \code{int n}, \code{iret} & interrupt, come back from it & IF = TF = 0 / restored \\ \hline
    \multicolumn{3}{|c|}{\cellcolor{brandOrange!20}\textbf{Strings}} \\ \hline
    \code{movsb}, \code{movsw} & \nql{ES:[DI]}{DS:[SI]}, then move SI and DI forward & none \\ \hline
    \code{rep} & repeat the instruction CX times & -- \\ \hline
    \code{cld}, \code{std} & DF = 0 (forward), DF = 1 & DF \\ \hline
\end{tabularx}
\end{center}

%================================================
\chapter{EMU8086 and NASM: differences in writing}
%================================================

The programs of Chapter 3 are written for EMU8086 (a syntax close to MASM), and the programs of Chapter 4 for NASM.
The processor is the same, but the two assemblers do not write some things the same way:

\begin{center}
\small
\begin{tabularx}{\linewidth}{|X|l|l|}
    \hline
    \ras{What we want} & \ras{EMU8086 / MASM} & \ras{NASM} \\ \hline
    address of a variable in a register & \code{lea dx, msg} or \code{mov dx, offset msg} & \code{mov dx, msg} \\ \hline
    value of a variable & \code{mov al, x} or \code{mov al, [x]} & \code{mov al, [x]} (brackets needed) \\ \hline
    giving the size & \code{mov byte ptr [bx], 5} & \code{mov byte [bx], 5} \\ \hline
    a hex number & \code{0FFh} & \code{0xFF} or \code{0FFh} \\ \hline
    a procedure & \code{p proc} \dots{} \code{p endp} & a simple label \code{p:} \\ \hline
    a constant & \code{N equ 10} & \code{N equ 10} \\ \hline
    an array of 10 zeros & \code{t db 10 dup(0)} & \code{t times 10 db 0} \\ \hline
    assembling & the \emph{Emulate} button & \code{nasm -f bin prog.asm -o prog.com} \\ \hline
\end{tabularx}
\end{center}

\begin{intabih}
    A good habit: always write the brackets \code{[ ]} when you mean the content of memory, even in EMU8086. Your program
    stays clear, and works with both assemblers with few changes. This is what we did in all the programs of this course.
\end{intabih}

%================================================
\chapter{Glossary}
%================================================

\begin{center}
\small\renewcommand{\arraystretch}{1.08}
\begin{tabular}{|l|l|}
    \hline
    \ras{English} & \ras{Français} \\ \hline
    computer & ordinateur \\ \hline
    processor, CPU & processeur, UCT \\ \hline
    central unit & unité centrale \\ \hline
    main memory & mémoire centrale \\ \hline
    volatile memory & mémoire volatile \\ \hline
    bus & bus \\ \hline
    address bus / data bus / control bus & bus d'adresses / de données / de contrôle \\ \hline
    register & registre \\ \hline
    program counter (PC) & compteur ordinal (CO) \\ \hline
    instruction register (IR) & registre d'instruction (RI) \\ \hline
    memory address register (MAR) & registre d'adresse mémoire (RAM) \\ \hline
    memory data register (MDR) & registre d'information mémoire (RIM) \\ \hline
    accumulator & accumulateur \\ \hline
    flag & indicateur, drapeau \\ \hline
    arithmetic logic unit (ALU) & unité arithmétique et logique (UAL) \\ \hline
    control unit & unité de commande (UC) \\ \hline
    decoder / sequencer & décodeur / séquenceur \\ \hline
    half adder / full adder & demi-additionneur / additionneur complet \\ \hline
    two's complement & complément à deux \\ \hline
    instruction set & jeu d'instructions \\ \hline
    opcode & code opération \\ \hline
    addressing mode & mode d'adressage \\ \hline
    effective address / physical address & adresse effective / adresse physique \\ \hline
    segment, offset & segment, déplacement \\ \hline
    instruction cycle & cycle d'instruction \\ \hline
    fetch, decode, execute & recherche, décodage, exécution \\ \hline
    assembly language / assembler & langage d'assemblage / assembleur \\ \hline
    stack & pile \\ \hline
    procedure, subroutine & procédure, sous-programme \\ \hline
    interrupt / interrupt vector table & interruption / table des vecteurs d'interruption \\ \hline
    I/O port & port d'entrée/sortie \\ \hline
    protected mode & mode protégé \\ \hline
    flat memory model & modèle mémoire plat \\ \hline
    pipeline / branch prediction & pipeline / prédiction de branchement \\ \hline
    core / thread & cœur / thread, fil d'exécution \\ \hline
    cache memory / cache line & mémoire cache / ligne de cache \\ \hline
    hit / miss & succès / défaut \\ \hline
    temporal and spatial locality & localité temporelle et spatiale \\ \hline
\end{tabular}
\end{center}

\section*{References}
\begin{enumerate}
    \item A. S. Tanenbaum, \emph{Structured Computer Organization}, Pearson, 2012.
    \item D. A. Patterson, J. L. Hennessy, \emph{Computer Organization and Design: The Hardware/Software Interface}, Morgan Kaufmann.
    \item P. Zanella, Y. Ligier, E. Lazard, \emph{Architecture et technologie des ordinateurs}, 6e édition, Dunod, 2018.
    \item B. B. Brey, \emph{The Intel Microprocessors: 8086/8088 \dots{} Core2 with 64-bit Extensions}, Pearson.
    \item EMU8086, Microprocessor Emulator.
    \item NASM, The Netwide Assembler: https://www.nasm.us/
    \item MASM, Microsoft Macro Assembler reference: \\ {\small https://learn.microsoft.com/en-us/cpp/assembler/masm/}
    \item GAS, The GNU Assembler: https://sourceware.org/binutils/docs/as/
\end{enumerate}
